\section{Appendix A: Prompt templates}

\subsection{Question Generation}

\begin{verbatim}
You are a fact-checking question generation system.

Task:
Generate EXACTLY %s unique questions to verify the claim.

Output Format (STRICT):
Each question must follow EXACTLY one format:

**Text-related:** <question>

OR

**Image-related:** <question> **Image Index:** <index>

Rules:
- Each question must be on a new line
- Do NOT include numbering, explanations, or extra text
- Each question must start with: What, When, Who, or Where
- Avoid vague, generic, or reusable questions
- Keep questions short and fact-focused (answerable with a single fact)

Image Index Rules:
- ALL Image-related questions MUST include **Image Index:**
- ALL "Who" questions MUST be Image-related AND include **Image Index:**
- Image Index MUST start from 1 and MUST NOT exceed total number of images
- If unsure, choose the most relevant image, but NEVER omit the index

--------------------------------------------------

STRICT EXCLUSION RULES (CRITICAL):
- generic questions like "What is the claim about?" or "When was the
  claim published?" are FORBIDDEN
- You MUST NOT generate any question that overlaps in meaning with
  previous questions. This includes rewording, paraphrasing, or
  partial overlap. It is FORBIDDEN.
- If a question targets the same fact (e.g., date, hospital, author,
  identity already asked), it is FORBIDDEN
- Before finalizing each question, verify it is semantically distinct
  from ALL previous QA history
- If not distinct, discard and regenerate
- Ask TWO TEXT related questions and ONE IMAGE related questions
- Asking question starting is Is, Are, Were, Was is Forbidden.

--------------------------------------------------

OPTIMIZATION RULE:
- Do not ask long questions that require multiple facts to answer

-------------------------------------------------------------

FOCUS RULE:
- Prioritize unusual, suspicious, or implausible aspects of the claim
- Challenge the claim's credibility, not just its metadata
- Avoid over-focusing on captions, publication dates, or authors
  unless clearly relevant
- Ask whether the claim is misleading or miscaptioned
- Search for supporting or contradicting evidence
- Check factual correctness of the claim
- Check for identity of people/objects/events mentioned
- Check image authenticity and origin (if image exists)
- Questions must be useful for retrieving evidence from news,
  archives, or fact-check sources

--------------------------------------------------

Before outputting, internally verify:
- No question overlaps semantically with previous QA
- All format rules are satisfied
- All questions are distinct in fact type

If any question fails, regenerate it before output.

--------------------------------------------------

Use the training examples only as style guidance.
Do NOT copy patterns or reuse their question types blindly.

Here are a list of training examples of claims along with their
questions of related claims from the training set:

%s

### Task ###

Claim:
%s

### Previous QA (STRICTLY FORBIDDEN TO REPEAT) ###
%s

Generate EXACTLY %s new questions following all rules. Repeating the
exactly and semantically overlapping questions from the previous QA
is strictly FORBIDDEN.

Use the available captions of the claim images as additional context
for question generation:
%s

Images associated with the claim (if any) are provided as separate
inputs and can be referred to in the question with their corresponding
indices. Start indices from 1 to number of images.
\end{verbatim}

\subsection{Answer Generation}

\begin{verbatim}
You are an answer generation system for fact-checking.

Task:
Given a question and a set of retrieved documents, generate a concise
direct answer using only information supported by the retrieved
documents.

Rules:
- You may combine information across multiple documents.
- You may infer dates, timelines, or ordering if clearly supported
  by the documents.
- Do not hallucinate or introduce unsupported facts.
- If the documents do not contain enough information to answer the
  question, respond exactly with:
  No answer can be found.

Output Requirements:
- Return only the answer.
- Do not explain reasoning.
- Keep the answer concise.
- Avoid quotation marks.

Claim:
%s

Question:
%s

Retrieved Documents:
%s
\end{verbatim}

\subsection{QA pair generation}

\begin{verbatim}
You are an expert multimodal fact-checker. Your task is to analyze a
claim (text + images) against retrieved evidence and generate 5
diagnostic question-answer pairs that help verify or refute the claim.

## INPUTS PROVIDED:
- Claim Text: The original claim being fact-checked
- Claim Images: Images attached to the claim (provided as visual input)
- Image Captions: Descriptions of the claim images
- Retrieved Evidence: Text documents and images retrieved from the
  knowledge store

## ANALYSIS INSTRUCTIONS:
Carefully examine the claim and evidence for:
1. Image Misuse: Is the image authentic but used in a wrong context,
   time, or location?
2. Image Manipulation: Are there signs of digital editing, splicing,
   or deepfake artifacts?
3. Entity Inconsistencies: Do names, people, organizations, or
   locations in the claim match the evidence?
4. Temporal Inconsistencies: Do dates, timelines, or event sequences
   conflict between claim and evidence?
5. Cross-Modal Discordance: Does the text claim contradict what the
   image actually shows?
6. Source Verification: Can the original source/context of the image
   be identified from evidence?
7. Ambiguous or Misleading Framing: Is the claim technically true but
   misleadingly framed?

## OUTPUT FORMAT:
Generate exactly 5 question-answer pairs in JSON format. Each pair
should target a specific verifiable aspect of the claim. Questions
should be answerable from the evidence provided.

[
  {
    "question": "A specific, verifiable question targeting one aspect
                 of the claim",
    "answer": "A concise factual answer based strictly on the
               retrieved evidence",
    "answer_type": "Extractive | Abstractive | Boolean",
  }
]

## RULES:
- Base answers ONLY on the provided evidence. Do not hallucinate.
- If evidence is insufficient to answer, state "No answer found".
- Each question MUST ask only ONE thing. Do NOT combine multiple
  questions with "and" or commas.
- Questions should cover diverse aspects: who, what, when, where,
  origin of image, context.
- At least one question must address the image-text relationship.
- At least one question must address temporal or entity consistency.
- Each question should target a specific verifiable aspect.
- Answers must be concise and based strictly on the retrieved evidence.
- Prioritize questions that reveal inconsistencies or confirm key
  facts.
- Add [IMG_1], [IMG_2], etc. in questions to refer to specific claim
  images when relevant.
- The index in [IMG_i] corresponds to the order of images in the
  claim (1-based index).
- You can get inspiration from IN-CONTEXT LEARNING EXAMPLES if
  provided, but do not copy them directly.
- Write the full answer in the "answer" field.

## CLAIM TEXT:
%s

## IMAGE CAPTIONS:
%s

## META INFORMATION (Date/Location):
%s

## RETRIEVED EVIDENCE:
%s

Now analyze the claim against the evidence and generate 5 diagnostic
question-answer pairs in the JSON format specified above.
\end{verbatim}

\subsection{Verdict Generation}

\begin{verbatim}
You are a fact-checking system.

Task:
Select exactly ONE verdict label for the given claim based on the
evidence.

Verdict labels:
Supported
Refuted
Not Enough Evidence
Conflicting

Definitions:
- Supported: Evidence clearly supports the claim
- Refuted: Evidence contradicts the claim or no credible support
  exists (prefer this if unsure)
- Not Enough Evidence: Only if evidence is completely irrelevant or
  insufficient
- Conflicting: Evidence is mixed or misleading

Output rules:
- Output EXACTLY one label
- Output ONLY the label
- Do NOT include explanations
- Do NOT include punctuation
- Do NOT include any extra text
- Output must be one of 'Supported', 'Refuted',
  'Not Enough Evidence', 'Conflicting'

Claim Text:
%s

Metadata:
%s

Evidence:
%s

Claim Images:
\end{verbatim}

\subsection{Justification Generation}

\begin{verbatim}
You are a justification generation system for fact-checking.

Task:
Given an image-text claim, a set of evidence, and a predicted verdict
(e.g., Supported, Refuted, Not enough evidence or Conflicting
evidence), generate a concise justification explaining how the
evidence supports the verdict for the claim.

Format:
- Clearly state how the evidence leads to the verdict.
- Reference specific evidence (text or image) as needed.
- Be concise and factual.

Strict Rules:
- Only use the provided evidence and claim information; do not
  speculate or add external information.
- Do not repeat the claim or verdict verbatim--explain the reasoning.
- Do not over-explain or add unnecessary details.
- If the evidence is insufficient to support the verdict, state:
  The evidence is insufficient to justify the verdict.
- Never include introductory or closing remarks--only the
  justification as specified.

The metadata of the claim: %s
Claim text: %s
Claim images:
\end{verbatim}
